% Clase 25 --- La constante dielectrica (y con ella n) depende de lambda (§7.3).
% La curva no pretende ser un ajuste: lo que importa es la TENDENCIA (n baja al
% crecer lambda) y que la variacion es chica pero suficiente para separar colores.
\begin{center}
\begin{tikzpicture}

  \begin{axis}[
    fisfig, width=10.5cm, height=5.6cm,
    xlabel={$\lambda$ \; (\si{\nano\metre})}, ylabel={$n = c/v = \sqrt{K_e}$},
    xmin=350, xmax=780, ymin=1.505, ymax=1.545,
    xtick={400,500,600,700}, ytick={1.51,1.52,1.53,1.54},
    yticklabel style={/pgf/number format/fixed,
                      /pgf/number format/precision=3},
  ]
    \addplot[figblue, smooth, samples=120, domain=370:770]
      {1.5053 + 4277/(x*x)};
    \node[figlblaux, anchor=south west] at (axis cs:455,1.5275)
      {vidrio (esquemático)};

    \draw[figaux] (axis cs:400,1.505) -- (axis cs:400,1.545);
    \draw[figaux] (axis cs:700,1.505) -- (axis cs:700,1.545);
    \node[figlblaux, violet, anchor=north west] at (axis cs:404,1.5435)
      {violeta: $n$ mayor $\Rightarrow$ $v$ menor};
    \node[figlblaux, red!70!black, anchor=south east] at (axis cs:697,1.5085)
      {rojo: $n$ menor $\Rightarrow$ $v$ mayor};
  \end{axis}

\end{tikzpicture}\\[0.5em]
{\small\itshape La razón física es que un dieléctrico tiene constante $K_e$
porque sus moléculas \textbf{se alinean} con el campo aplicado, y en una onda ese
campo está oscilando. Con un campo estático las moléculas tienen todo el tiempo
del mundo para acomodarse; con un campo que cambia $10^{15}$ veces por segundo,
para cuando alcanzan a orientarse el campo ya se dio vuelta. Por eso la $K_e$
efectiva \textbf{no es} la de la electrostática del curso, y depende de la
frecuencia: $K_e = K_e(\lambda)$ y, con ella, $n = \sqrt{K_e}$. La variación es
chica ---acá, dos centésimas entre un extremo y otro del visible--- pero alcanza
para que el violeta viaje más lento que el rojo dentro del vidrio, que es
exactamente la \textbf{dispersión cromática} de la clase siguiente: la razón por
la que un prisma abre la luz blanca en colores y por la que toda tabla de índices
de refracción viene con una $\lambda$ especificada al lado (la de referencia
suele ser la línea amarilla del sodio, \SI{589}{\nano\metre}).}
\end{center}
